\pdfoutput=1
\documentclass[11pt,a4paper]{article}

% ---- Encoding & Fonts (IACR recommended) ----
\usepackage[T1]{fontenc}
\usepackage[utf8]{inputenc}
\usepackage{lmodern}

% ---- Math ----
\usepackage{amsmath, amssymb, amsthm}
\usepackage{mathtools}

% ---- Page Layout ----
\usepackage[top=1.2in,bottom=1.2in,left=1.2in,right=1.2in]{geometry}

% ---- Hyperlinks ----
\usepackage{hyperref}
\hypersetup{
    colorlinks=true,
    linkcolor=blue,
    filecolor=magenta,
    urlcolor=cyan,
    citecolor=red,
    pdftitle={Star-Mesh: A Hybrid Post-Quantum, Decentralized Forward-Secret Messaging Protocol},
    pdfauthor={Samin Yasar},
    pdfkeywords={Post-Quantum Cryptography, Secure Messaging, Hybrid Ratchet, Decentralized Networks,
                 Kademlia DHT, ML-KEM, ML-DSA},
}

% ---- Author affiliations ----
\usepackage{authblk}

% ---- ORCID icon ----
\usepackage{orcidlink}

% ---- Algorithms ----
\usepackage{algorithm}
\usepackage{algorithmic}

% ---- Figures ----
\usepackage{graphicx}

% ---- References ----
\usepackage{cite}

% ---- Code listings ----
\usepackage{listings}
\usepackage{xcolor}

% ---- Tables ----
\usepackage{booktabs}

% ---- Typography ----
\usepackage{microtype}

% ---- URLs in bibliography ----
\usepackage{url}

% ============================================================
% Theorem environments
% ============================================================
\newtheorem{theorem}{Theorem}[section]
\newtheorem{securityclaim}[theorem]{Security Claim}
\newtheorem{definition}[theorem]{Definition}
\newtheorem{lemma}[theorem]{Lemma}
\newtheorem{corollary}[theorem]{Corollary}
\newtheorem{proposition}[theorem]{Proposition}
\theoremstyle{remark}
\newtheorem{remark}[theorem]{Remark}
\theoremstyle{plain}
\newtheorem{assumption}[theorem]{Assumption}

% ============================================================
% Cryptographic macros
% ============================================================
\newcommand{\negl}{\mathsf{negl}}
\newcommand{\Adv}{\mathbf{Adv}}
\newcommand{\PGame}{\mathbf{Game}}
\newcommand{\calA}{\mathcal{A}}
\newcommand{\calB}{\mathcal{B}}
\newcommand{\calS}{\mathcal{S}}
\newcommand{\calK}{\mathcal{K}}
\newcommand{\calO}{\mathcal{O}}
\newcommand{\bits}{\{0,1\}}
\newcommand{\getsr}{\stackrel{\$}{\leftarrow}}
\newcommand{\MLKEM}{\mathsf{ML\text{-}KEM}}
\newcommand{\MLDSA}{\mathsf{ML\text{-}DSA}}
\newcommand{\HKDF}{\mathsf{HKDF}}
\newcommand{\DHop}{\mathsf{DH}}
\newcommand{\Encaps}{\mathsf{Encaps}}
\newcommand{\Decaps}{\mathsf{Decaps}}
\newcommand{\Sign}{\mathsf{Sign}}
\newcommand{\Verify}{\mathsf{Verify}}
\newcommand{\KeyGen}{\mathsf{KeyGen}}
% Fixed-width byte-length encoding
\newcommand{\flen}[1]{\mathsf{len}_{32}(#1)}
% Hybrid game shorthand
\newcommand{\Hgame}[1]{\mathbf{H}_{#1}}

% ============================================================
% Title block
% ============================================================
\title{\textbf{Star-Mesh: A Hybrid Post-Quantum, Decentralized\\
               Forward-Secret Messaging Protocol}}

\author{Samin Yasar~\orcidlink{0009-0004-5292-523X}}
\affil{Independent Researcher \\ Dhaka, Bangladesh \\
       \texttt{research@samin-yasar.dev}}
\date{August 13, 2026}

\begin{document}

\maketitle

% ============================================================
% Abstract
% ============================================================
\begin{abstract}
The transition to post-quantum cryptography necessitates a fundamental reevaluation of secure
messaging protocols, particularly those reliant on decentralized infrastructure.  In this paper
we present \emph{Star-Mesh}, a fully decentralized, peer-to-peer messaging protocol that
provides robust \emph{hybrid} post-quantum confidentiality, forward secrecy, and
post-compromise security.  Unlike centralized architectures, Star-Mesh operates over a Kademlia
Distributed Hash Table (DHT) for key distribution and leverages a mixnet-based onion-routing
overlay for metadata obfuscation.

We make the following contributions.  While Star-Mesh builds upon established components such
as X3DH, the Double Ratchet, Kademlia routing, and Loopix-style mixnets, the integration is
non-trivial: the large size of post-quantum keys and signatures directly conflicts with the
strict constant-packet-size discipline of Loopix/Sphinx mixnets (which pad all packets to a
fixed MTU to prevent traffic analysis), meaning that simply layering PQ-KEM epochs onto such a
network forces packet fragmentation and degrades anonymity.  Star-Mesh resolves this tension
through a systematic NIST Category~3 parameter selection that maximizes post-quantum security
margins without exceeding standard mixnet packetization thresholds (see Section~\ref{sec:paramselection}).
The further substantive contribution is a formal treatment of freshness and recovery semantics
in a decentralized hybrid post-quantum messaging setting: we unify the analysis of RNG-compromise
freshness, post-compromise recovery, and the distinction between session-level and one-time
pre-key PQ-FS.  This is the point at which the protocol is differentiated from prior hybrid
designs, and it directly addresses the primary concern that a hybrid or semi-static pre-key
design can weaken forward secrecy semantics.
(1)~We design a \emph{hybrid} key-exchange mechanism combining classical X25519
Diffie--Hellman with ML-KEM-768~\cite{nist_fips203} under a concatenation combiner bound by
HKDF, using one-time post-quantum pre-keys to guarantee initial forward secrecy against a
quantum-capable adversary.
(2)~We provide a formal specification of a hybrid post-quantum Double Ratchet state machine
incorporating ephemeral ML-KEM exchanges, and give game-based security definitions for
Hybrid Post-Quantum Forward Secrecy (PQ-FS) and Post-Compromise Security (PCS) that explicitly
integrate RNG compromise into the freshness condition.
(3)~We sketch game-based security arguments: PQ-FS is argued to follow from the IND-CCA2
security of ML-KEM-768, PCS restoration is argued after one complete ratchet round-trip (with
explicit caveats on out-of-order delivery and the skipped-key cache), and authentication
security is argued under EUF-CMA security of ML-DSA-65~\cite{nist_fips204}.
(4)~We provide a bandwidth and performance analysis showing that the PQ overhead amortises to
approximately $22.7 + \delta_{\text{DHT}}$ bytes per message at ratchet interval $\Delta = 50$,
where $\delta_{\text{DHT}}$ denotes the DHT routing overhead.
\end{abstract}

\paragraph{Keywords:}
Post-Quantum Cryptography; Secure Messaging; Hybrid Ratchet; Decentralized Networks;
Kademlia DHT; ML-KEM; ML-DSA; Forward Secrecy; Post-Compromise Security.

\tableofcontents
\newpage

% ============================================================
\section{Introduction}
% ============================================================

The imminent threat of Cryptographically Relevant Quantum Computers (CRQCs) has accelerated the
standardization of post-quantum cryptographic (PQC) algorithms by NIST, notably ML-KEM
\cite{nist_fips203} and ML-DSA \cite{nist_fips204}.  While traditional secure messaging
protocols such as Signal have begun integrating PQC \cite{signal_pqxdh} to defend against
``Harvest Now, Decrypt Later'' (HNDL) attacks \cite{mosca2018cybersecurity}, they rely
predominantly on centralized servers for key distribution and message routing.

This centralization creates a single point of failure for both availability and metadata
privacy: a server compromise can reveal the social graph, timing information, and message volume
of all users, even when message contents are encrypted.  Decentralized alternatives such as
Session \cite{session_protocol} and Briar \cite{briar_design} mitigate this architectural risk
through onion-routing networks, but they remain committed exclusively to classical
elliptic-curve primitives (X25519), leaving them vulnerable to HNDL attacks.

\paragraph{Our Approach.}
Star-Mesh synthesizes three independently mature technologies: (1)~NIST-standardized
post-quantum cryptography \cite{nist_fips203,nist_fips204}, (2)~Kademlia DHT \cite{kademlia}
for decentralized key distribution, and (3)~mixnet-style onion routing \cite{chaum_mix,loopix}
for metadata resistance.  The cryptographic core is a \emph{hybrid} PQ-X3DH handshake followed
by a hybrid PQ Double Ratchet, building on the classical Double Ratchet analysis of Cohn-Gordon
et al.\ \cite{cohngordon2016} and the Signal X3DH specification \cite{marlinspike_x3dh}.

\paragraph{Hybrid Security Rationale.}
Star-Mesh is a \emph{hybrid} post-quantum protocol: every handshake and every DH ratchet step
combines classical X25519 with ML-KEM-768.  The term ``hybrid'' is precise and intentional: the
protocol is \emph{not} fully post-quantum, because X25519 remains a required component.
Security holds if \emph{either} component is secure, which provides both defense-in-depth
against a classical adversary (should an early ML-KEM break be found) and protection against a
CRQC-equipped adversary (via ML-KEM, whose security reduces to Module-LWE
\cite{lyubashevsky2010mlwe}).

Formally, the hybrid shared secret is computed as
\[
  SS_{\text{hybrid}} = \texttt{0xFF} \;\|\; SS_{cl} \;\|\; SS_{PQ1} \;\|\; SS_{PQ2}
\]
passed through HKDF-SHA3-256 with the ciphertexts $ct_1, ct_2$ and the classical ephemeral public key $EK_{pk,A}$ bound into the KDF input string $info$. Under this model, following the analysis of Giacon et al.\ \cite{giacon2018kem}, the concatenation combiner ensures that the derived key is pseudorandom and IND-CCA2 secure if either $SS_{cl}$ or the combination of $SS_{PQ1}$ and $SS_{PQ2}$ remains secure under the stated assumptions; binding the transcript (ephemeral key and ciphertexts) satisfies the theorem's precondition, defending against related-ciphertext/re-encapsulation attacks. HKDF is modeled as a PRF in all proof sketches (see Assumption~\ref{asm:hkdf} in Section~\ref{sec:secmodel}). Our security claims rely specifically on the concatenation-combiner setting analyzed in \cite{giacon2018kem}.

\paragraph{Trust Assumptions on Infrastructure.}
Star-Mesh is \emph{decentralized in deployment}, not trust-free.  The DHT nodes collectively
control availability of key bundles; a Sybil attack can eclipse a user's mailbox.  The mixnet
requires that fewer than a fraction $f$ of mix nodes are compromised (the Loopix assumption
\cite{loopix}).  These operational assumptions are stated explicitly in
Section~\ref{sec:assumptions}.

\paragraph{Contributions.}
\begin{enumerate}
  \item \textbf{Formal freshness model:} A unified, RNG-integrated game-based treatment of
        post-quantum forward secrecy and post-compromise recovery, explicitly capturing the
        distinction between session-level and one-time pre-key PQ-FS under mixed $PQ\_SPK$ and
        $PQ\_OTPK$ usage.
  \item \textbf{Authenticated pre-key design:} A hybrid PQ-X3DH bundle format with signed
        batch hashing of all OTPK and PQ-OTPK material, closing the key-substitution and
        identity-misbinding vectors from the general UKS attack class highlighted in the
        literature \cite{blake1997unknown}.
  \item \textbf{Protocol design:} A fully specified hybrid PQ-X3DH handshake and PQ Double
        Ratchet for asynchronous P2P messaging over a Kademlia DHT, integrating the classical
        X25519 and ML-KEM components in a single decentralized architecture.  This design
        resolves the non-trivial conflict between post-quantum payload sizes and the strict
        constant-packet-size requirement of Loopix/Sphinx mixnets through an explicit
        Category~3 parameter analysis (Section~\ref{sec:paramselection}).
  \item \textbf{Security analysis:} Game-based security arguments for PQ-FS and PCS against
        IND-CCA2 of ML-KEM-768, and for authentication against EUF-CMA of ML-DSA-65.  All
        hybrid arguments are spelled out with named intermediate games.
  \item \textbf{Metadata resistance under PQC overhead:} Analysis of the mixnet overlay
        achieving sender-receiver unlinkability under standard mixnet assumptions \cite{loopix},
        with an explicit argument that Category~3 parameter sizes preserve the constant-packet-size
        invariant on which Loopix's anonymity guarantees depend.
  \item \textbf{Performance analysis:} Concrete bandwidth cost accounting with an explicit
        $\delta_{\text{DHT}}$ term bounding DHT routing overhead.
  \item \textbf{Prototype implementations:} Tested Python and Rust proof-of-concept
        implementations demonstrating OKM convergence and ratchet state transitions
        (Section~\ref{sec:poc}).
\end{enumerate}

\paragraph{Scope and Proof Status.}
This work is intended as a research-oriented protocol design and security-analysis preprint
for public dissemination on a preprint server.  The game-based arguments in
Sections~\ref{sec:security} should be interpreted as structured proof sketches rather than
fully mechanized or reduction-tight proofs in a composable framework.  In particular, we do
not claim a complete UC-style treatment of the protocol \cite{uc_framework}, nor a full
machine-verified analysis of the ratchet state machine.  Our goal is to provide a precise
cryptographic design together with explicit assumptions, adversarial capabilities, and
security-reduction intuition sufficient to support further peer review, implementation, and
formalization work.

\paragraph{Research Positioning.}
Star-Mesh should be viewed as a protocol-design and analysis contribution rather than a
production-ready messaging system.  The present manuscript focuses on architectural synthesis,
adversarial modeling, and hybrid post-quantum ratcheting semantics in a preprint setting,
where the goal is to surface a precise design and a clear security model for peer review.
A production deployment would additionally require extensive implementation hardening,
empirical evaluation, interoperability testing, formal verification, and operational
abuse-resistance analysis.

\paragraph{Paper Organization.}
Section~\ref{sec:prelim} introduces notation and primitives.
Section~\ref{sec:protocol} specifies the Star-Mesh protocol.
Section~\ref{sec:secmodel} gives formal security definitions and assumptions.
Section~\ref{sec:security} presents the security claims and proof sketches.
Section~\ref{sec:metadata} analyses metadata resistance.
Section~\ref{sec:performance} provides a bandwidth and performance analysis.
Section~\ref{sec:related} surveys related work and Section~\ref{sec:conclusion} concludes.

% ============================================================
\section{Preliminaries}
\label{sec:prelim}
% ============================================================

\subsection{Notation}

We write $\lambda$ for the security parameter.  A function $f(\lambda)$ is \emph{negligible}
(written $f \in \negl(\lambda)$) if for every polynomial $p$, $f(\lambda) < 1/p(\lambda)$ for
all sufficiently large $\lambda$.  We write $x \getsr S$ for sampling $x$ uniformly at random
from a finite set $S$, and $x \getsr \mathcal{D}$ for sampling from a distribution
$\mathcal{D}$.  We denote by $\|$ the concatenation of byte strings.

We write $\flen{s}$ for the \emph{fixed-width, 4-byte big-endian} encoding of the byte length
of string $s$.  Using a fixed-width encoding prevents length-extension ambiguity when multiple
variable-length fields are concatenated.

For key derivation, $\HKDF(ikm, salt, info, \ell)$ denotes HKDF-SHA3-256 with input key
material $ikm$, optional $salt$, context $info$, and output length $\ell$ bytes.  SHA3-256 is
the underlying hash function; it is \emph{not} modeled as a random oracle in this work; rather
HKDF is modeled as a two-stage extract-then-expand PRF following the HKDF construction in
RFC~5869 \cite{rfc5869} and its formal analysis by Krawczyk (CRYPTO 2010)
\cite{krawczyk2010hkdf} (see Assumption~\ref{asm:hkdf} in
Section~\ref{sec:secmodel}).

\subsection{Key Encapsulation Mechanisms}

\begin{definition}[KEM]
A Key Encapsulation Mechanism $\Pi = (\KeyGen, \Encaps, \Decaps)$ consists of three algorithms:
\begin{itemize}
  \item $\KeyGen(1^\lambda) \to (pk, sk)$: Generates a public/secret key pair.
  \item $\Encaps(pk) \to (ct, ss)$: Encapsulates, producing ciphertext $ct$ and shared secret $ss$.
  \item $\Decaps(sk, ct) \to ss$: Decapsulates, recovering $ss$ (or $\bot$ on failure).
\end{itemize}
\end{definition}

\begin{definition}[IND-CCA2 KEM Security]
Let $\Pi = (\KeyGen, \Encaps, \Decaps)$ be a KEM.  The advantage of adversary $\calA$ in the
IND-CCA2 game is:
\[
  \Adv^{\mathrm{ind\text{-}cca2}}_{\Pi}(\calA) \coloneqq \left| \Pr\left[
  \begin{array}{l}
    (pk, sk) \getsr \KeyGen(1^\lambda); \\
    (ct^*, ss_0) \getsr \Encaps(pk); \\
    ss_1 \getsr \bits^{|ss_0|}; \\
    b \getsr \bits; \\
    b' \getsr \calA^{\Decaps(sk,\cdot)}(pk, ct^*, ss_b) : \\
    \quad b' = b
  \end{array}
  \right] - \frac{1}{2} \right|
\]
where $\calA$ may not query $\Decaps(sk, ct^*)$.  We say $\Pi$ is IND-CCA2 secure if
$\Adv^{\mathrm{ind\text{-}cca2}}_{\Pi}(\calA) \in \negl(\lambda)$ for all PPT $\calA$.
\end{definition}

\subsection{Digital Signature Schemes}

\begin{definition}[EUF-CMA Security]
A signature scheme $\Sigma = (\KeyGen, \Sign, \Verify)$ is Existentially Unforgeable under
Chosen-Message Attacks (EUF-CMA) if for all PPT adversaries $\calA$:
\[
  \Adv^{\mathrm{euf\text{-}cma}}_{\Sigma}(\calA) \coloneqq \Pr\left[
  \begin{array}{l}
    (vk, sk) \getsr \KeyGen(1^\lambda); \\
    (m^*, \sigma^*) \getsr \calA^{\Sign(sk,\cdot)}(vk) : \\
    \quad \Verify(vk, m^*, \sigma^*) = 1 \\
    \quad \wedge\ m^* \notin \mathcal{Q}
  \end{array}
  \right] \in \negl(\lambda)
\]
where $\mathcal{Q}$ is the set of queries made to the signing oracle.
\end{definition}

\subsection{Cryptographic Primitives in Star-Mesh}

Star-Mesh instantiates the following standardized primitives:

\begin{itemize}
  \item \textbf{Identity Signatures:} ML-DSA-65 (Dilithium3) \cite{nist_fips204}, a
        lattice-based signature scheme whose security reduces to the hardness of Module-LWE and
        Module-SIS.
  \item \textbf{Key Encapsulation:} ML-KEM-768 (Kyber768) \cite{nist_fips203}, an IND-CCA2
        KEM whose security reduces to Module-LWE with parameters $n = 256$, $k = 3$.  Public
        key size: 1184 bytes; ciphertext size: 1088 bytes; shared secret: 32 bytes.
  \item \textbf{Classical DH:} X25519 \cite{rfc7748}, a Diffie--Hellman function over
        Curve25519.  Key and output size: 32 bytes each.  X25519 is vulnerable to a CRQC; it
        is retained for hybrid security against classical adversaries only.
  \item \textbf{Key Derivation:} HKDF-SHA3-256 \cite{rfc5869} for all key derivation steps
        (see Assumption~\ref{asm:hkdf}); BLAKE3 \cite{blake3} for message key derivation within
        the symmetric ratchet.
  \item \textbf{AEAD Encryption:} AES-256-GCM \cite{nist_aesgcm} with 96-bit nonces and
        128-bit tags.
\end{itemize}

\subsection{Parameter Selection and Security Categories}
\label{sec:paramselection}

To establish concrete security levels and estimate protocol overhead, Star-Mesh instantiates FIPS 203 and FIPS 204 parameters aligned with NIST Security Category~3 (providing key-search security equivalent to AES-192). Specifically, we select ML-KEM-768 for post-quantum key encapsulation and ML-DSA-65 for identity signatures. This selection represents a deliberate balance of security, communication bandwidth, and computational cost under decentralized mixnet constraints:

\begin{itemize}
  \item \textbf{Comparison with Category 1 (ML-KEM-512 and ML-DSA-44):} While Category~1 primitives provide smaller public key sizes (e.g., 800~B for ML-KEM-512) and marginally faster execution times, they only offer AES-128 equivalent security. We opt for Category~3 to align with standard high-security targets (AES-192/256) and to provide a wider margin of safety against future algebraic/structural advances in lattice attacks.
  \item \textbf{Comparison with Category 5 (ML-KEM-1024 and ML-DSA-87):} Category~5 provides the highest standardized security (AES-256 equivalent) but introduces substantial bandwidth penalties. ML-DSA-87 signatures are over 4.5~KB, and ML-KEM-1024 public keys and ciphertexts are 1.5~KB each. In a decentralized mixnet architecture like Loopix/Sphinx, messages are routed in constant-sized packets (MTUs) to prevent traffic analysis. Introducing massive signatures would force packet fragmentation or require significantly larger packet padding sizes, severely degrading overall mixnet throughput and increasing network propagation delays.
\end{itemize}

Thus, Category~3 (ML-KEM-768 and ML-DSA-65) is selected as the default parameter set for Star-Mesh, maximizing security without exceeding standard mixnet MTU packetization thresholds.

% ============================================================
\section{The Star-Mesh Protocol}
\label{sec:protocol}
% ============================================================

\subsection{Overview and Trust Model}
\label{sec:trust}

Star-Mesh is a decentralized protocol deployed without central servers; all key-distribution
operations are performed over a Kademlia DHT \cite{kademlia}.  Users are identified by their
ML-DSA-65 public key $IK^{DSA}$.  Out-of-band verification (e.g., key fingerprint comparison,
QR codes) is assumed for initial identity binding, following standard practice in secure
messaging \cite{marlinspike_x3dh}.

\textbf{DHT trust assumption.}  The DHT stores and serves Pre-Key Bundles but is \emph{not}
trusted for their authenticity: all bundles are signed under the owner's $IK^{DSA}_{sk}$ and
verified by recipients.  However, DHT nodes collectively control availability; a Sybil attack
can eclipse a user's mailbox.  We assume the DHT is Sybil-resistant as an operational
assumption (see Section~\ref{sec:assumptions}).

\textbf{Mixnet trust assumption.}  Message delivery is routed through a mixnet overlay.  We
assume at most a fraction $f < 1/2$ of mix nodes are adversarially controlled, following the
Loopix threat model \cite{loopix}.

\subsection{Key Generation and Identity}

Each node derives its cryptographic material from a 128-bit master seed
$\mathsf{seed} \getsr \bits^{128}$ via deterministic key generation.  A node's \emph{Identity
Bundle} comprises:
\begin{itemize}
  \item $IK^{DSA} = (IK^{DSA}_{pk}, IK^{DSA}_{sk}) \getsr \MLDSA.\KeyGen(1^\lambda)$: the
        long-term identity signing keypair.
  \item $IK^{DH} = (IK^{DH}_{pk}, IK^{DH}_{sk}) \getsr X25519.\KeyGen$: the long-term
        classical DH keypair.
  \item $SPK = (SPK_{pk}, SPK_{sk}) \getsr X25519.\KeyGen$: a semi-static signed pre-key,
        rotated periodically (e.g., weekly).
  \item $PQ\_SPK = (PQ\_SPK_{pk}, PQ\_SPK_{sk}) \getsr \MLKEM.\KeyGen(1^\lambda)$: a
        semi-static ML-KEM pre-key, rotated after a bounded number of uses.
  \item $\{OTPK_i\}_{i=1}^{n}$: a batch of $n$ one-time X25519 pre-keys.
  \item $\{PQ\_OTPK_j\}_{j=1}^{m}$: a batch of $m$ one-time ML-KEM pre-keys.  Consumption of
        at least one $PQ\_OTPK$ is \emph{mandatory} when available to guarantee initial PQ
        forward secrecy (see Section~\ref{sec:fsremark}).  When the batch is exhausted, only
        $PQ\_SPK$ is used; in that case, the session achieves \emph{session-level} PQ-FS (the
        pre-key is ephemeral per session) rather than \emph{one-time pre-key PQ-FS} (the
        pre-key is used exactly once, then zeroized).  See Remark~\ref{rem:spkonly} for the
        precise distinction.
\end{itemize}

\subsection{Pre-Key Bundle Publication}
\label{sec:bundle}

Node $A$ publishes the following signed bundle to the DHT, keyed by $H(IK^{DSA}_{pk,A})$:
\[
  \mathsf{Bundle}_A = \bigl(IK^{DSA}_{pk}, IK^{DH}_{pk}, SPK_{pk}, PQ\_SPK_{pk},
                             \{OTPK_{pk,i}\}_{i}, \{PQ\_OTPK_{pk,j}\}_{j}, \sigma\bigr)
\]
The signature $\sigma$ is computed over \emph{all} public pre-key material to prevent
substitution of individual pre-keys by a DHT adversary:
\begin{align*}
  \mathsf{BundlePayload} &=
        SPK_{pk} \;\|\; PQ\_SPK_{pk} \\
        &\quad\|\; \flen{\bigoplus_{i} OTPK_{pk,i}}
             \;\|\; H\!\left(\bigoplus_{i} OTPK_{pk,i}\right) \\
        &\quad\|\; \flen{\bigoplus_{j} PQ\_OTPK_{pk,j}}
             \;\|\; H\!\left(\bigoplus_{j} PQ\_OTPK_{pk,j}\right) \\[4pt]
  \sigma &= \MLDSA.\Sign(IK^{DSA}_{sk},\; \mathsf{BundlePayload})
\end{align*}
where $H$ is SHA3-256 and $\bigoplus$ denotes ordered concatenation of all keys in the batch.
Signing the hashes of the full OTPK and PQ-OTPK batches (rather than individual keys) binds
the entire batch to $\sigma$ without inflating the signed payload size; an adversary cannot
substitute any individual OTPK or PQ-OTPK without invalidating $\sigma$.

Any node fetching $\mathsf{Bundle}_A$ verifies $\sigma$ against the reconstructed
$\mathsf{BundlePayload}$ before use.

\begin{remark}
\label{rem:spkonly}
\textbf{Session-level vs.\ one-time pre-key PQ-FS.}  When a PQ-OTPK is consumed, its secret
is zeroized immediately after a single decapsulation; subsequent sessions receive a distinct
one-time pre-key.  This gives \emph{one-time pre-key PQ-FS} for the handshake: even revealing
the long-term $IK^{DH}_{sk}$ and $PQ\_SPK_{sk}$ after the fact does not expose the session's
shared secret, because $PQ\_OTPK_{sk}$ is gone.

When the OTPK batch is exhausted, the initiator encapsulates only against $PQ\_SPK$.  Because
$PQ\_SPK$ is reused across multiple sessions before rotation, an adversary that obtains
$PQ\_SPK_{sk}$ can retroactively recover the shared secret of all sessions that used it.  We
term this \emph{session-level} PQ-FS: forward secrecy holds at the granularity of $PQ\_SPK$
rotation epochs, not per session.  Implementations should replenish OTPK batches proactively to
minimize exposure in this mode.
\end{remark}

\subsection{Hybrid PQ-X3DH Handshake}
\label{sec:x3dh}

Alice (initiator) fetches Bob's bundle from the DHT, verifies $\sigma_B$ over the full
$\mathsf{BundlePayload}_B$, and executes the following handshake.  The design extends Signal's
X3DH \cite{marlinspike_x3dh} with parallel ML-KEM encapsulations.

\paragraph{Step 1 --- Classical X3DH Component.}
Alice generates a fresh ephemeral keypair $EK_A \getsr X25519.\KeyGen$ and computes:
\begin{align*}
  DH_1 &= \DHop(IK^{DH}_{sk,A},\; SPK_{pk,B}) \\
  DH_2 &= \DHop(EK_{sk,A},\; IK^{DH}_{pk,B}) \\
  DH_3 &= \DHop(EK_{sk,A},\; SPK_{pk,B}) \\
  DH_4 &= \DHop(EK_{sk,A},\; OTPK_{pk,B}) \quad \text{(if available; omitted otherwise)}
\end{align*}
The classical shared secret is $SS_{cl} = DH_1 \| DH_2 \| DH_3 \| DH_4$.

\paragraph{Step 2 --- Post-Quantum KEM Component.}
Alice encapsulates against Bob's ML-KEM pre-keys:
\begin{align*}
  (ct_1,\; SS_{PQ1}) &\getsr \MLKEM.\Encaps(PQ\_SPK_{pk,B}) \\
  (ct_2,\; SS_{PQ2}) &\getsr \MLKEM.\Encaps(PQ\_OTPK_{pk,B}) \quad \text{(if available)}
\end{align*}

\paragraph{Step 3 --- Cryptographic Binding and Key Derivation.}
The Associated Data binds the handshake to both parties' post-quantum signing identities,
preventing identity misbinding and Unknown Key-Share (UKS)-class attacks (originally
demonstrated for STS) \cite{blake1997unknown}:
\[
  AD = \texttt{``StarMesh''} \| \flen{IK^{DSA}_{pk,A}} \| IK^{DSA}_{pk,A}
                           \| \flen{IK^{DSA}_{pk,B}} \| IK^{DSA}_{pk,B}
\]
To satisfy the IND-CCA2 security precondition of the concatenation/split-key combiner proved by Giacon et al.~\cite{giacon2018kem}, the KEM ciphertexts and the classical ephemeral public key are bound into the HKDF input. The binder string $info$ is constructed by appending the length-prefixed transcript parameters to $AD$:
\[
  info = AD \| \flen{EK_{pk,A}} \| EK_{pk,A} \| \flen{ct_1} \| ct_1 \| \flen{ct_2} \| ct_2
\]
The hybrid input key material is formed by combining classical and post-quantum secrets with a
domain-separation prefix byte $\texttt{0xFF}$ (following \cite{bindel2018hybrid}):
\[
  SS_{\text{hybrid}} = \texttt{0xFF} \| SS_{cl} \| SS_{PQ1} \| SS_{PQ2}
\]
The 64-byte Output Keying Material is then:
\[
  OKM = \HKDF\!\left(ikm = SS_{\text{hybrid}},\; salt = \mathbf{0}^{32},\;
                     info = info,\; \ell = 64\right)
\]
\textbf{Rationale for $salt = \mathbf{0}^{32}$ at initialization.}  RFC~5869
\cite{rfc5869} states that when no truly independent salt is available (as is the case at the
start of a fresh handshake where no prior shared state exists), a zero salt is acceptable: the
HKDF extract step still provides entropy concentration from the IKM.  In contrast, subsequent
ratchet steps use $salt = RK$ (the current Root Key), which is a secret, independent input that
provides additional forward-secrecy benefit.  The different salt usage in the handshake vs.\ the
ratchet is therefore intentional and semantically justified.

The first 32 bytes of $OKM$ initialize the Root Key $RK$; the second 32 bytes initialize the
sending Chain Key $CK_s$.

\paragraph{Handshake Message.}
Alice transmits to Bob:
\[
  M_0 = \bigl(IK^{DSA}_{pk,A},\; IK^{DH}_{pk,A},\; EK_{pk,A},\; ct_1,\; ct_2,\;
              \text{identifier of consumed pre-keys}\bigr)
\]
Including $IK^{DH}_{pk,A}$ directly in $M_0$ (rather than requiring Bob to resolve it via a
separate DHT bundle lookup) preserves the protocol's asynchronous, offline-completable handshake
property: Bob can derive the full session key from $M_0$ alone, with no live dependency on DHT
availability at the moment he reads his mailbox.

\noindent\textbf{Implicit authentication in $M_0$.}  $M_0$ is not explicitly signed.  However,
the computed $SS_{cl}$ includes $DH_1 = \DHop(IK^{DH}_{sk,A}, SPK_{pk,B})$, which is computable
only by the holder of $IK^{DH}_{sk,A}$.  Bob's successful derivation of the same $OKM$ thus
implicitly authenticates Alice's identity, analogously to the implicit authentication property
of standard X3DH \cite{marlinspike_x3dh}.  A replay of $M_0$ with a substituted $EK_{pk,A}$
would fail because the $DH_2/DH_3/DH_4$ components are computed from $EK_{sk,A}$, which the
replayer does not possess.

Bob uses the $IK^{DH}_{pk,A}$ included in $M_0$ together with his own $SPK_{sk,B}$ to recompute
$DH_1$, decapsulates $ct_1, ct_2$ using the corresponding secret keys, recomputes
$SS_{\text{hybrid}}$, and derives the same $OKM$, completing the handshake.

\subsection{Ephemeral PQ Double Ratchet}
\label{sec:ratchet}

Following handshake completion, both parties enter the PQ Double Ratchet.  The state machine
extends the classical Signal Double Ratchet \cite{marlinspike_ratchet} by interleaving ML-KEM
epochs with X25519 DH ratchet steps.

\subsubsection{Ratchet State}

Each node maintains a session state $\mathcal{S}$:
\begin{itemize}
  \item $RK \in \bits^{32}$: Root Key.
  \item $CK_s, CK_r \in \bits^{32}$: Sending and Receiving Chain Keys.
  \item $DH^{loc}_{sk}, DH^{rem}_{pk}$: Current local X25519 secret key and remote public key.
  \item $N_s, N_r \in \mathbb{N}$: Send and receive message counters.
  \item $PQ^{loc}_{sk}, PQ^{loc}_{pk}$: Current local ephemeral ML-KEM keypair (may be $\bot$).
  \item $PQ^{pending}_{ct}$: Pending ML-KEM ciphertext received from peer, awaiting decapsulation.
  \item $\mathcal{K}_{skipped}$: Cache of skipped message keys, bounded at capacity 1000.
  \item $\mathsf{pq\_step\_counter} \in \{0, \ldots, \Delta\}$: Countdown to the next mandatory
        PQ ratchet step, with $\Delta = 50$.
\end{itemize}

\subsubsection{Symmetric Ratchet (Message Key Derivation)}

To encrypt message $i$ in the current sending chain:
\begin{align*}
  MK_i \| CK_s' &= \mathsf{BLAKE3\text{-}KDF}(CK_s, \texttt{``StarMesh-MK''}) \\
  CK_s &\leftarrow CK_s' \quad (\text{then zeroize previous } CK_s) \\
  N_s &\leftarrow N_s + 1
\end{align*}
The message is encrypted as
$c = \mathsf{AES\text{-}256\text{-}GCM}_{MK_i}(m;\; \mathsf{nonce} = N_s^{32} \| \texttt{0}^{64})$,
where $N_s^{32}$ is the 4-byte big-endian encoding of the 32-bit send counter $N_s$, and
$\texttt{0}^{64}$ is 8 zero bytes, together forming a 12-byte (96-bit) nonce as required by
AES-256-GCM.  The 32-bit counter field is byte-aligned and caps a single chain at
$2^{32}-1 \approx 4.3 \times 10^9$ messages before nonce reuse; in practice, chains are
replaced by DH ratchet steps and PQ ratchet steps well before this bound, but the limit
is stated here as a deliberate design choice rather than an implementation artifact.
$MK_i$ is zeroized immediately after use. We instantiate the symmetric ratchet using BLAKE3-KDF rather than HKDF-SHA256 due to the high frequency of symmetric updates (triggered per message). BLAKE3's tree-hashing structure yields significantly lower latency and power consumption on resource-constrained devices, while natively supporting secure domain separation.

\subsubsection{DH Ratchet Step}

When Alice sends after receiving a message from Bob (i.e., when $DH^{rem}_{pk}$ has been
updated), a DH ratchet step is triggered:
\[
  DH_{\mathrm{new}} = \DHop(DH^{loc}_{sk},\; DH^{rem}_{pk})
\]
\[
  RK,\; CK_s = \HKDF\!\left(ikm = DH_{\mathrm{new}},\; salt = RK,\;
                             info = \texttt{``StarMesh-RK''},\; \ell = 64\right)
\]
Alice generates a new $DH^{loc}_{sk}$ and zeroizes the previous one.  The new $DH^{loc}_{pk}$
is attached to the next outgoing message.

\subsubsection{PQ Ratchet Step}
\label{sec:pqratchet}

The PQ ratchet step is triggered when $\mathsf{pq\_step\_counter} = 0$ or on the first message
of a new session.  It proceeds as follows:

\begin{enumerate}
  \item \textbf{Key Generation:} Alice generates a fresh ML-KEM keypair
        $(PQ^{loc}_{sk}, PQ^{loc}_{pk}) \getsr \MLKEM.\KeyGen(1^\lambda)$ and attaches
        $PQ^{loc}_{pk}$ to the next outgoing message.
  \item \textbf{Encapsulation (Bob's side):} Upon receiving $PQ^{loc}_{pk}$, Bob computes
        $(PQ^{pending}_{ct}, SS_{PQ}) \getsr \MLKEM.\Encaps(PQ^{loc}_{pk})$ and attaches
        $PQ^{pending}_{ct}$ to his reply.
  \item \textbf{Root Chain Update (Alice's side):} Upon receiving $PQ^{pending}_{ct}$, Alice
        computes $SS_{PQ} = \MLKEM.\Decaps(PQ^{loc}_{sk}, PQ^{pending}_{ct})$ and updates:
        \[
          RK,\; CK_s = \HKDF\!\left(ikm = SS_{PQ},\; salt = RK,\;
                                     info = \texttt{``StarMesh-PQ-RK''},\; \ell = 64\right)
        \]
        Alice immediately zeroizes $PQ^{loc}_{sk}$.
  \item \textbf{Counter Reset:} $\mathsf{pq\_step\_counter} \leftarrow \Delta$.
\end{enumerate}

\begin{remark}[Justification of $\Delta = 50$ and CRQC Exposure Window]
\label{rem:delta}
The parameter $\Delta$ governs the maximal interval between PQ ratchet steps.  Over this
interval, session confidentiality against a CRQC relies solely on the X25519 DH ratchet.

\textbf{Worst-case exposure bound.}  In the worst case, if a session transmits $\Delta$ messages
before the PQ ratchet fires, a CRQC that breaks all $\Delta$ DH ratchet steps can compute the
current chain key $CK$ from the captured traffic.  From $CK$, the adversary can derive all
$\Delta$ message keys in the current symmetric chain.  However, message keys from \emph{prior}
DH ratchet epochs are \emph{not} recoverable (assuming the deletion schedule is enforced): they
depend on previous DH outputs that are no longer derivable from the current DH state, and the
corresponding chain keys have been zeroized.  Thus, the CRQC exposure window for a single DH
epoch is at most $\Delta$ messages, bounded to the epoch between two consecutive DH ratchet steps.

Choosing $\Delta = 50$ amortises the 1184-byte ML-KEM public key cost over 50 messages while
bounding the CRQC exposure window to 50 ciphertexts per epoch.  Operators with higher security
requirements may reduce $\Delta$ at the cost of increased bandwidth.
\end{remark}

\subsubsection{Out-of-Order Message Handling}

If a message arrives with sequence number $N > N_r$, the receiving node derives and caches the
skipped message keys $\{MK_{N_r}, \ldots, MK_{N-1}\}$ in $\mathcal{K}_{skipped}$.  The cache
is bounded at 1000 entries; if full, the oldest cached key is zeroized and dropped (the
corresponding message is permanently undecryptable).  When a delayed message arrives, the node
retrieves its key from $\mathcal{K}_{skipped}$, decrypts, and immediately zeroizes the key.

\subsubsection{Deletion Schedule}

The following deletion schedule is mandatory and forms the basis of the forward secrecy and
PCS proof sketches.  Its enforcement is an \emph{operational assumption} (see
Section~\ref{sec:assumptions}); the protocol logic specifies the schedule but cannot itself
verify that an implementation honours it.

\begin{itemize}
  \item \textbf{Message Keys ($MK_i$):} Zeroized immediately after a single AES-GCM
        encryption or decryption operation.
  \item \textbf{Chain Keys ($CK_s$, $CK_r$):} Zeroized immediately upon derivation of the
        subsequent chain key.
  \item \textbf{Ephemeral X25519 Secrets ($DH^{loc}_{sk}$):} Zeroized immediately upon
        generation of the next local DH keypair.
  \item \textbf{Ephemeral ML-KEM Secret ($PQ^{loc}_{sk}$):} Zeroized immediately after
        decapsulation of the peer's $PQ^{pending}_{ct}$, i.e., within one full round-trip.
  \item \textbf{One-Time Pre-Key Secrets:} Zeroized immediately after use in the handshake.
  \item \textbf{Skipped Keys:} Zeroized either upon successful use or upon eviction from the
        bounded cache.
\end{itemize}

\subsection{Prototype Implementations}
\label{sec:poc}

To validate the correctness of the cryptographic state machine described in this section, we
developed two proof-of-concept implementations available in the open-source repository
accompanying this paper\footnote{\url{https://github.com/Samin-yasar/star-mesh} (Apache-2.0 License).}:

\begin{enumerate}
  \item \textbf{Python Prototype} (\texttt{poc/python/poc.py}): A dependency-free Python~3
        implementation that simulates the commutative key exchange and the state transitions of
        the hybrid Double Ratchet.  DH operations are modeled via modular exponentiation over a
        standard 1024-bit MODP prime (RFC~2409 Group~2) to preserve the algebraic commutativity
        property of Diffie--Hellman, and ML-KEM-768 is simulated while preserving the
        IND-CCA2 encapsulate/decapsulate API contract.  The Python prototype successfully
        verifies OKM convergence between initiator and responder.

  \item \textbf{Rust Prototype} (\texttt{poc/rust/src/main.rs}): A high-fidelity implementation
        using production-ready RustCrypto libraries:
        \begin{itemize}
          \item \texttt{ml-kem}~(v0.3) for FIPS~203 ML-KEM-768 \cite{nist_fips203}.
          \item \texttt{x25519-dalek}~(v2.0) for RFC~7748 X25519 \cite{rfc7748}.
          \item \texttt{blake3}~(v1.0) and \texttt{hkdf}~(v0.12) / \texttt{sha2}~(v0.10) for
                BLAKE3-KDF and HKDF-SHA256.
          \item \texttt{zeroize}~(v1.0) for memory shredding, directly enforcing the deletion
                schedule specified above.
        \end{itemize}
\end{enumerate}

Both prototypes execute the full hybrid PQ-X3DH handshake, verify OKM convergence, advance the
symmetric and post-quantum ratchet steps, and confirm root-key convergence after a PQ ratchet
round-trip, consistent with the PCS healing mechanism argued in Section~\ref{sec:pcsproof}.

% ============================================================
\section{Security Model}
\label{sec:secmodel}
% ============================================================

\subsection{Cryptographic Assumptions}
\label{sec:assumptions}

We state all assumptions explicitly; proof sketches reference only assumptions listed here.

\begin{assumption}[Module-LWE Hardness]
\label{asm:mlwe}
The Module-LWE problem with parameters $n = 256$, $k = 3$ (ML-KEM-768 parameters) is hard for
all QPT adversaries.  Under this assumption, ML-KEM-768 is IND-CCA2 secure.
\end{assumption}

\begin{assumption}[HKDF as PRF]
\label{asm:hkdf}
HKDF-SHA3-256 is modeled as a pseudorandom function (PRF) in the standard model, following the
two-stage (extract-then-expand) analysis of Krawczyk \cite{krawczyk2010hkdf}.  All proof sketches
invoking the pseudorandomness of derived keys rely on this assumption.  We do \emph{not} model
SHA3-256 as a random oracle; the PRF modeling is sufficient for our reductions.  In particular,
for this work we model the full interface $\HKDF(ikm, salt, info, \ell)$ as PRF-secure under the
instantiated salt choices used by the protocol (including $salt = \mathbf{0}^{32}$ at
initialization and $salt = RK$ in ratchet updates).
\end{assumption}

\begin{assumption}[DHT Availability and Sybil Resistance]
\label{asm:dht}
The Kademlia DHT is available and Sybil-resistant in the sense that no adversary can eclipse
a target user's mailbox.  This is an operational assumption; it is not proven by the protocol
logic.
\end{assumption}

\begin{assumption}[Mixnet Honesty Fraction]
\label{asm:mixnet}
At most a fraction $f < 1/2$ of mix nodes are adversarially controlled, following the Loopix
threat model \cite{loopix}.  This quantifies the trust requirement on the routing overlay.
\end{assumption}

\begin{assumption}[Deletion Schedule Enforcement]
\label{asm:deletion}
The deletion schedule specified in Section~\ref{sec:ratchet} is enforced by the implementation.
In particular, message keys, chain keys, and ephemeral secret keys are zeroized at the times
specified.  This is an operational assumption.
\end{assumption}

\begin{assumption}[RNG Integrity after $t_{rec}$]
\label{asm:rng}
The node's random number generator is not compromised by the adversary for any sampling event
after time $t_{rec}$.  This models a recovered operating environment (clean OS RNG,
side-channel-free execution) following a prior compromise event.  Foundational theory for
extracting usable randomness from imperfect, seed-dependent sources is given by
Dodis et al.\ \cite{dodis2013entropic}; our $\mathcal{O}_{\mathsf{RevealRNG}}$/recovery-time
abstraction is an explicit operational modelling assumption of this work.  If the RNG is
compromised during the PQ ratchet
key generation after $t_{rec}$, PCS is information-theoretically unachievable.
\end{assumption}

\subsection{Execution Environment and Adversary}

We model execution using a multi-party protocol execution game adapted from
\cite{cohngordon2016,dowling2020ratchet}.  Oracle naming in this paper is normalized to a
single notation ($\calO^{Reveal}$, $\calO^{RevealLTK}$, etc.); it is intended to be
semantically aligned with their state-reveal and long-term-key compromise queries rather than a
verbatim reproduction of prior oracle labels.  Let $\mathcal{P} = \{P_1, \ldots, P_n\}$ be a set
of parties.  Each party $P_i$ may hold multiple sessions; we denote by $\pi^s_i$ the $s$-th
session of party $P_i$.

We evaluate Star-Mesh against a \textbf{quantum-active adversary} $\calA_{\mathrm{quant}}$,
modeled as a non-uniform quantum polynomial-time (QPT) algorithm equipped with a CRQC.
$\calA_{\mathrm{quant}}$ has the following capabilities.

\begin{definition}[Adversarial Oracles]
\label{def:oracles}
The adversary $\calA_{\mathrm{quant}}$ may adaptively query the following oracles:
\begin{itemize}
  \item $\calO^{Send}(i, s, m)$: Delivers, injects, modifies, or drops message $m$ to/from
        session $\pi^s_i$.  $\calA_{\mathrm{quant}}$ has full network control.
  \item $\calO^{Reveal}(i, s)$: Returns the current ratchet state $\mathcal{S}$ of session
        $\pi^s_i$ (simulating device compromise).
  \item $\calO^{RevealLTK}(i)$: Returns the long-term secret keys of party $P_i$.
  \item $\calO^{RevealRNG}(i, t)$: Reveals or controls the output of $P_i$'s random number
        generator at time $t$.  This models a compromised OS RNG, a side-channel attack, or
        any mechanism by which the adversary learns or influences fresh randomness
        \cite{dodis2013entropic}.
  \item $\calO^{Test}(i, s, e)$: Returns either the real epoch-$e$ message key $MK_e$ or a
        uniformly random key; the adversary's goal is to distinguish.
\end{itemize}
\end{definition}

\begin{definition}[Fresh Session Epoch]
\label{def:fresh}
An epoch $e$ of session $\pi^s_i$ is \textbf{fresh} if no $\calO^{Test}$ query has been made
on it, and the adversary's combined query history does not allow derivation of $MK_e$ from
revealed material.  Formally, epoch $e$ is \emph{not fresh} if any of the following hold:
\begin{enumerate}
  \item $\calO^{Reveal}(i, s)$ was queried after the last DH ratchet step preceding epoch $e$,
        \emph{and} $\calO^{Reveal}(j, s')$ was queried for the corresponding peer session
        $\pi^{s'}_j$.
  \item $\calO^{RevealLTK}(i)$ was queried \emph{and} all one-time pre-key secrets of
        \emph{Bob's} bundle used in session $\pi^s_i$ were previously revealed via
        $\calO^{Reveal}$.
  \item $\calO^{RevealRNG}(i, t)$ was queried for any time $t$ at which a fresh ephemeral
        secret (ephemeral DH key, ephemeral ML-KEM keypair, or one-time pre-key) used to derive
        $MK_e$ was sampled.
\end{enumerate}
\end{definition}

\begin{remark}
Condition~3 integrates RNG compromise into the freshness definition, replacing the ad-hoc RNG
boundary used in prior formulations.  It ensures that the PCS game (Definition~\ref{def:pcs})
does not need to re-define the RNG boundary separately: any session whose PQ ratchet keypair
was sampled during a compromised-RNG window is simply not fresh for the affected epochs.
\end{remark}

\subsection{Post-Quantum Forward Secrecy (PQ-FS) Game}

\begin{definition}[PQ-FS Security]
\label{def:pqfs}
Consider the following game $\PGame^{\mathrm{PQ\text{-}FS}}_{\calA}(\lambda)$:
\begin{enumerate}
  \item The challenger runs all honest party key generation and initializes the DHT.
  \item $\calA_{\mathrm{quant}}$ adaptively issues $\calO^{Send}$, $\calO^{Reveal}$,
        $\calO^{RevealLTK}$, and $\calO^{RevealRNG}$ queries.
  \item $\calA_{\mathrm{quant}}$ selects a target session $\pi^s_i$ and a target epoch $e$
        that is fresh per Definition~\ref{def:fresh}.  It additionally requires that $e$
        precedes the time of any $\calO^{Reveal}$ query on $\pi^s_i$.
  \item The challenger samples $b \getsr \bits$.  If $b = 0$, it returns $MK_e$; if $b = 1$,
        it returns $r \getsr \bits^{256}$.
  \item $\calA_{\mathrm{quant}}$ outputs a guess $b'$.
\end{enumerate}
The advantage is $\Adv^{\mathrm{PQ\text{-}FS}}_{\mathrm{SM}}(\calA) \coloneqq
|\Pr[b' = b] - \tfrac{1}{2}|$.  The protocol achieves \textbf{PQ-FS} if this advantage is
negligible for all QPT $\calA_{\mathrm{quant}}$.
\end{definition}

\subsection{Post-Compromise Security (PCS) Game}

\begin{definition}[PCS Security]
\label{def:pcs}
Consider the following game $\PGame^{\mathrm{PCS}}_{\calA}(\lambda)$:
\begin{enumerate}
  \item The challenger initializes all sessions.  $\calA_{\mathrm{quant}}$ may issue arbitrary
        oracle queries.
  \item $\calA_{\mathrm{quant}}$ issues $\calO^{Reveal}(i, s)$ at time $t_{comp}$, obtaining
        the full ratchet state.
  \item $\calO^{RevealRNG}(i, t)$ may be queried at any time $t \leq t_{rec}$ for some
        recovery time $t_{rec} \geq t_{comp}$.  After $t_{rec}$, by Assumption~\ref{asm:rng},
        the RNG is clean; freshness Condition~3 (Definition~\ref{def:fresh}) therefore does not
        apply to keys sampled after $t_{rec}$.
  \item After time $t_{rec}$, at least one complete PQ ratchet round-trip is executed: Alice
        samples $(PQ^{loc}_{sk}, PQ^{loc}_{pk})$ freshly after $t_{rec}$; sends $PQ^{loc}_{pk}$
        to Bob; Bob responds with $PQ^{pending}_{ct}$; Alice decapsulates and zeroizes
        $PQ^{loc}_{sk}$.
  \item $\calA_{\mathrm{quant}}$ selects a target epoch $e'$ in session $\pi^s_i$ satisfying
        (a)~$e'$ occurs after the round-trip of step~4 completes; (b)~the chain key for epoch
        $e'$ is derived \emph{after} the round-trip (not from a pre-compromise chain key); and
        (c)~the message key $MK_{e'}$ is not in $\mathcal{K}_{skipped}$ from before $t_{rec}$.
        The challenger responds with a Test oracle call.
  \item $\calA_{\mathrm{quant}}$ outputs a guess $b'$.
\end{enumerate}
The protocol achieves \textbf{PCS} if $\Adv^{\mathrm{PCS}}_{\mathrm{SM}}(\calA) \in
\negl(\lambda)$ for all QPT $\calA_{\mathrm{quant}}$.
\end{definition}

% ============================================================
\section{Security Analysis}
\label{sec:security}
% ============================================================

\subsection{Authentication: MITM Resistance}
\label{sec:fsremark}

\begin{securityclaim}[Identity Binding]
\label{thm:auth}
If there exists a QPT adversary $\calA_{\mathrm{quant}}$ that forges a valid Pre-Key Bundle or
executes a successful Man-in-the-Middle attack on the PQ-X3DH handshake with advantage
$\epsilon$, then there exists a QPT reduction $\calB$ that breaks EUF-CMA security of ML-DSA-65
with advantage $\epsilon' \geq \epsilon - \negl(\lambda)$.
\end{securityclaim}

\begin{proof}[Proof Sketch]
Let $\calB$ be a reduction interacting with an EUF-CMA challenger $\mathcal{C}$ for ML-DSA-65.
$\calB$ receives a challenge verification key $vk^*$ and embeds it as $IK^{DSA}_{pk,A}$ in the
simulated P2P network.  Signing queries from $\calA_{\mathrm{quant}}$ are forwarded to
$\mathcal{C}$'s signing oracle.

Suppose $\calA_{\mathrm{quant}}$ succeeds in either: (a)~publishing a malicious bundle
$\mathsf{Bundle}'_A$ with a valid signature under $IK^{DSA}_{pk,A}$, or (b)~altering the $AD$
parameter (which includes $IK^{DSA}_{pk,A}$ by construction) such that the resulting $OKM$ does
not correspond to an honest handshake, while convincing an honest participant that it does.

In case~(a): the bundle signature now covers the full $\mathsf{BundlePayload}$ including hashes
of OTPK and PQ-OTPK batches (Section~\ref{sec:bundle}).  If $\calA_{\mathrm{quant}}$ publishes
$\mathsf{Bundle}'_A$ with any substituted pre-key (whether $SPK$, $PQ\_SPK$, any $OTPK_i$, or
any $PQ\_OTPK_j$), the corresponding batch hash in $\mathsf{BundlePayload}$ changes, so
$\sigma'$ must be a fresh signature on a message not in $\mathcal{Q}$.  $\calB$ extracts
$(m', \sigma')$ and submits it to $\mathcal{C}$, winning the EUF-CMA game.

In case~(b): modifying $AD$ would require $\calA_{\mathrm{quant}}$ to produce a bundle verified
under $IK^{DSA}_{pk,A}$ with a different pre-key payload, which is again a signature forgery.

Therefore $\Adv^{\mathrm{euf\text{-}cma}}_{\MLDSA}(\calB) \geq \epsilon - \negl(\lambda)$,
where the negligible term accounts for the probability that $\calA_{\mathrm{quant}}$'s forgery
message was already in $\mathcal{Q}$.
\end{proof}

\subsection{Post-Quantum Forward Secrecy}

\begin{securityclaim}[PQ-FS]
\label{thm:pqfs}
Under Assumptions~\ref{asm:mlwe} and \ref{asm:hkdf}, if there exists a QPT adversary
$\calA_{\mathrm{quant}}$ winning $\PGame^{\mathrm{PQ\text{-}FS}}_{\calA}(\lambda)$
(Definition~\ref{def:pqfs}) with advantage $\epsilon$, then there exists a QPT reduction
$\calB$ that breaks IND-CCA2 security of ML-KEM-768 with advantage
$\epsilon' \geq \epsilon/2 - \negl(\lambda)$.
\end{securityclaim}

\begin{proof}[Proof Sketch]
We analyze security via a sequence of hybrid games.

\medskip\noindent
\textbf{Game $\Hgame{0}$:} The real PQ-FS game $\PGame^{\mathrm{PQ\text{-}FS}}_{\calA}(\lambda)$.
$\calA_{\mathrm{quant}}$'s advantage in $\Hgame{0}$ is $\epsilon$ by hypothesis.

\medskip\noindent
\textbf{Game $\Hgame{1}$:} Identical to $\Hgame{0}$, except that the challenger, upon choosing
Bob's bundle for the target session $\pi^s_i$, embeds the IND-CCA2 challenge public key $pk^*$
as $PQ\_OTPK_{pk,B}$.  The corresponding ciphertext $ct^*$ (from the IND-CCA2 challenger)
replaces $ct_2$ in Alice's handshake message $M_0$.

\textit{Transition $\Hgame{0} \to \Hgame{1}$:} From $\calA_{\mathrm{quant}}$'s view, this
substitution is syntactically identical to the real protocol: $pk^*$ is a valid ML-KEM public
key, and $ct^*$ is a valid encapsulation thereof.  No oracle is modified.  The two games are
therefore \emph{perfectly indistinguishable}:
$|\Pr[\calA\text{ wins }\Hgame{1}] - \Pr[\calA\text{ wins }\Hgame{0}]| = 0$.

\medskip\noindent
\textbf{Game $\Hgame{2}$:} Identical to $\Hgame{1}$, except that $SS_{PQ2}$ is replaced by a
uniformly random string $u \getsr \bits^{32}$, independently of $ct^*$.  This corresponds to
the IND-CCA2 challenger providing the ``random'' branch ($ss_1$) rather than the real shared
secret.

\textit{Transition $\Hgame{1} \to \Hgame{2}$:}  In $\Hgame{1}$, $SS_{PQ2}$ is the real
decapsulation of $ct^*$ under $pk^*$.  In $\Hgame{2}$, $SS_{PQ2}$ is random.  Suppose
$\calA_{\mathrm{quant}}$ distinguishes $\Hgame{1}$ from $\Hgame{2}$ with advantage $\delta$.
Then the reduction $\calB$, who receives $(pk^*, ct^*, ss_b)$ from the IND-CCA2 challenger,
runs $\calA_{\mathrm{quant}}$ in the $\Hgame{1}$ simulation but uses $ss_b$ in place of
$SS_{PQ2}$.  $\calB$ outputs $b' = 0$ if $\calA_{\mathrm{quant}}$ wins, and $b' = 1$ otherwise.
The IND-CCA2 decapsulation oracle handles all $\calA_{\mathrm{quant}}$'s decapsulation queries
on $ct \neq ct^*$; the query $ct = ct^*$ is not permitted by the IND-CCA2 rules (and
$\calA_{\mathrm{quant}}$'s PQ-FS freshness constraint prevents Reveal queries that would expose
$PQ\_OTPK_{sk}$ after decapsulation).  Thus $\Adv^{\mathrm{ind\text{-}cca2}}_{\MLKEM}(\calB)
\geq \delta$.

\medskip\noindent
\textbf{Game $\Hgame{3}$:} In $\Hgame{2}$, $SS_{\text{hybrid}} = \texttt{0xFF} \| SS_{cl} \|
SS_{PQ1} \| u$.  By Assumption~\ref{asm:hkdf} (HKDF as PRF), the output
$OKM = \HKDF(SS_{\text{hybrid}}, \mathbf{0}^{32}, AD, 64)$ is pseudorandom, hence $MK_e$
(derived from $OKM$) is pseudorandom and indistinguishable from a fresh random key.
$\calA_{\mathrm{quant}}$'s advantage in $\Hgame{3}$ is therefore $\negl(\lambda)$.

\textit{Transition $\Hgame{2} \to \Hgame{3}$:}  The distinguishing advantage between
$\Hgame{2}$ and $\Hgame{3}$ is at most $\Adv^{\mathrm{PRF}}_{\HKDF}(\cdot) \in \negl(\lambda)$
by Assumption~\ref{asm:hkdf}.

\medskip\noindent
\textbf{Bounding $\epsilon'$:}  The reduction $\calB$ does not know a priori which of Bob's
ML-KEM pre-keys ($PQ\_SPK$ or $PQ\_OTPK$) will be the one encapsulated by Alice.  $\calB$
guesses the target key by embedding $pk^*$ as $PQ\_OTPK_{pk,B}$; this guess is correct with
probability at least $1/2$ (when both $PQ\_SPK$ and $PQ\_OTPK$ are available, $\calB$ guesses
which encapsulation $\calA_{\mathrm{quant}}$ relies on).  The $1/2$ loss is the standard
single-target guessing penalty for a reduction that must select one pre-key branch among
multiple valid candidates; it is a technique limitation of the proof, not an apparent protocol
weakness.  A multi-instance or hybrid reduction that simultaneously handles the full set of
candidate ML-KEM pre-keys removes this loss and yields the tight bound up to negligible terms:
\[
  \Adv^{\mathrm{ind\text{-}cca2}}_{\MLKEM}(\calB)
    \;\geq\; \frac{1}{2}\bigl(\epsilon - \negl(\lambda)\bigr)
    \;=\; \frac{\epsilon}{2} - \negl(\lambda).
\]

Combining the hybrids:
$\epsilon = |\Pr[\calA\text{ wins }\Hgame{0}] - \Pr[\calA\text{ wins }\Hgame{3}]|
           \leq \delta + \negl(\lambda)$,
so $\delta \geq \epsilon - \negl(\lambda)$, and
$\epsilon' \geq \epsilon/2 - \negl(\lambda)$.
\end{proof}

\begin{remark}[Role of the Classical Component]
Following the analysis of Giacon et al.\ \cite{giacon2018kem} (concatenation combiner with PRF
KDF, instantiated with HKDF per Assumption~\ref{asm:hkdf} where the ciphertexts and ephemeral keys are bound to the KDF input), if X25519 remains
computationally secure (no CRQC), then $SS_{cl}$ alone is expected to make $SS_{\text{hybrid}}$
pseudorandom under the stated assumptions, providing standard (classical) forward secrecy.
Binding the ciphertexts $ct_1, ct_2$ into $info$ satisfies the precondition of \cite{giacon2018kem},
guarding against related-ciphertext/re-encapsulation attacks.
The ML-KEM component provides additional PQ-FS against a CRQC.  The two components are
\emph{complementary}: security is argued to hold if either is secure under the modeled
assumptions.
\end{remark}

\subsection{Post-Compromise Security}
\label{sec:pcsproof}

\begin{securityclaim}[PCS Healing Bound]
\label{thm:pcs}
Under Assumptions~\ref{asm:mlwe}, \ref{asm:hkdf}, \ref{asm:deletion}, and \ref{asm:rng},
suppose $\calA_{\mathrm{quant}}$ wins $\PGame^{\mathrm{PCS}}_{\calA}(\lambda)$
(Definition~\ref{def:pcs}) with advantage $\epsilon$.  Then there exists a QPT reduction
$\calB$ that breaks IND-CCA2 security of ML-KEM-768 with advantage
$\epsilon' \geq \epsilon/2 - \negl(\lambda)$.
\end{securityclaim}

\begin{proof}[Proof Sketch]
We again proceed by a sequence of hybrid games.

\medskip\noindent
\textbf{Game $\Hgame{0}$:} The real PCS game.  $\calA_{\mathrm{quant}}$'s advantage is
$\epsilon$.

\medskip\noindent
\textbf{Game $\Hgame{1}$:} Identical to $\Hgame{0}$, except that the challenger embeds the
IND-CCA2 challenge public key $pk^*$ as $PQ^{loc}_{pk}$ (Alice's fresh post-$t_{rec}$ ML-KEM
public key) and the challenge ciphertext $ct^*$ as $PQ^{pending}_{ct}$.  Since $PQ^{loc}_{pk}$
is generated after $t_{rec}$ (clean RNG, Assumption~\ref{asm:rng}), no prior state snapshot
obtained via $\calO^{Reveal}$ contains $PQ^{loc}_{sk}$.

\textit{Transition $\Hgame{0} \to \Hgame{1}$:}  Perfectly indistinguishable (same argument as
Security Claim~\ref{thm:pqfs}).

\medskip\noindent
\textbf{Game $\Hgame{2}$:} Replace $SS_{PQ} = \MLKEM.\Decaps(PQ^{loc}_{sk}, ct^*)$ with
$u \getsr \bits^{32}$.  This corresponds to the IND-CCA2 challenger's ``random'' branch.

\textit{Transition $\Hgame{1} \to \Hgame{2}$:}  The deletion schedule (Assumption~\ref{asm:deletion})
mandates that $PQ^{loc}_{sk}$ is zeroized immediately after decapsulation.  Any subsequent
$\calO^{Reveal}$ query yields only post-zeroization state, containing no information about
$SS_{PQ}$.  The reduction $\calB$ is constructed identically to Security Claim~\ref{thm:pqfs}: it
uses $\calA_{\mathrm{quant}}$'s distinguishing advantage on target epoch $e'$ (which satisfies
conditions (a)--(c) of Definition~\ref{def:pcs}) to break the IND-CCA2 challenge.

\medskip\noindent
\textbf{Game $\Hgame{3}$:} By Assumption~\ref{asm:hkdf}, the new root key
$RK' = \HKDF(u, RK, \texttt{``StarMesh-PQ-RK''}, 64)[\text{bytes }0\text{--}31]$ is
pseudorandom, and all subsequent chain keys $CK'$ and message keys derived from $RK'$ are
pseudorandom.  $\calA_{\mathrm{quant}}$'s advantage is $\negl(\lambda)$.

\medskip\noindent
\textbf{Bounding $\epsilon'$:}  By the same $1/2$-guessing argument as
Security Claim~\ref{thm:pqfs} (here $\calB$ embeds $pk^*$ as $PQ^{loc}_{pk}$, which is the
unique ML-KEM key involved in the post-recovery round-trip; the $1/2$ factor arises because
$\calB$ must guess which direction carries the relevant encapsulation in a symmetric protocol):
\[
  \epsilon' \;\geq\; \frac{\epsilon}{2} - \negl(\lambda).
\]
\end{proof}

\begin{corollary}[PCS Healing]
\label{cor:pcs}
Under the conditions of Security Claim~\ref{thm:pcs}, confidentiality of all epochs $e'$
satisfying conditions \emph{(a)}, \emph{(b)}, and \emph{(c)} of Definition~\ref{def:pcs} is
restored after one complete PQ ratchet round-trip, under the IND-CCA2 security of ML-KEM-768.

In particular, note the following limitations.  First, out-of-order delivery: if Bob's
$PQ^{pending}_{ct}$ reply is delayed, the round-trip may not complete within a bounded window;
healing is delayed accordingly.  Second, skipped-key cache: message keys cached in
$\mathcal{K}_{skipped}$ from epochs before $t_{rec}$ are not covered by this corollary (those
message keys are already exposed by the prior compromise).  Third, PQ ratchet counter: if
$\mathsf{pq\_step\_counter} > 0$ at $t_{comp}$, Alice may not send a fresh $PQ^{loc}_{pk}$
until the counter expires (at most $\Delta$ additional messages), delaying the start of the
round-trip.
\end{corollary}

% ============================================================
\section{Metadata Resistance}
\label{sec:metadata}
% ============================================================

Star-Mesh achieves metadata resistance through two complementary mechanisms.

\paragraph{DHT Mailbox Indirection.}
Each node maintains a \emph{mailbox} in the Kademlia DHT keyed by a pseudonymous address
$\mathsf{addr}_i = H(IK^{DSA}_{pk,i} \| \mathsf{epoch})$, rotated periodically.  Senders
deposit encrypted messages at the mailbox address; recipients perform constant-rate polling.
This severs the direct IP-to-recipient link exploitable by a passive observer of DHT traffic.

\paragraph{Onion-Routing Overlay.}
Message delivery and mailbox retrieval are routed through a libp2p-based onion-routing overlay
following the Loopix \cite{loopix} architecture.  Each message is layered with Sphinx-format
encryption \cite{sphinx} over a randomly sampled sequence of mix nodes.  Dummy cover traffic is
generated at a constant rate $\lambda_{cover}$ to prevent traffic analysis via timing or volume.

\begin{proposition}[Sender-Receiver Unlinkability]
Under Assumption~\ref{asm:mixnet} (at most $f < 1/2$ of mix nodes are adversarially controlled)
and assuming constant-rate cover traffic, the probability that $\calA_{\mathrm{quant}}$ links a
sending event to a receiving event is at most $\frac{f}{1-f} \cdot \frac{1}{n_{\mathrm{mix}}} +
\negl(\lambda)$, where $n_{\mathrm{mix}}$ is the number of honest mix nodes on the path.
\end{proposition}

\noindent The bound above is a \emph{derived approximation used in this work}, in the style of
the corruption-fraction threat model underlying Loopix \cite{loopix}.  We do not claim that
Section~4.3 of \cite{loopix} states this closed-form expression; that section reports
empirical end-to-end anonymity evaluation under Poisson mixing and constant-rate cover traffic.
Our use of $\frac{f}{1-f} \cdot \frac{1}{n_{\mathrm{mix}}}$ is therefore an explicit
modelling assumption under Assumption~\ref{asm:mixnet}, included to parameterise system-level
risk in Star-Mesh.  We make no new claims about mixnet security beyond this approximation.
Star-Mesh's contribution is the integration of this overlay with the hybrid post-quantum
ratchet layer.

% ============================================================
\section{Bandwidth and Performance Analysis}
\label{sec:performance}
% ============================================================

A critical challenge in adopting PQC over P2P networks is the substantially increased size of
ciphertexts and keys relative to classical primitives.  Table~\ref{tab:primitives} summarises
the relevant ML-KEM-768 and X25519 sizes.

\begin{table}[ht]
\centering
\caption{Primitive sizes relevant to Star-Mesh.}
\label{tab:primitives}
\begin{tabular}{@{}llll@{}}
\toprule
\textbf{Primitive} & \textbf{Public Key} & \textbf{Ciphertext / DH Output} & \textbf{Shared Secret} \\
\midrule
X25519    & 32 B    & 32 B              & 32 B \\
ML-KEM-768 & 1184 B & 1088 B            & 32 B \\
ML-DSA-65 & 1952 B  & N/A (sig: 3{,}309 B) & N/A  \\
\bottomrule
\end{tabular}
\end{table}

Table~\ref{tab:overhead} compares the per-connection overhead of a classical Signal-style
protocol against Star-Mesh.

\begin{table}[ht]
\centering
\caption{Bandwidth overhead comparison.  Classical Signal handshake size is derived from the
X3DH specification \cite{marlinspike_x3dh}: two 32-byte X25519 public keys (IK and EK), one
32-byte SPK, one 32-byte OTPK, and protocol framing, totalling approximately 150 bytes (118
bytes without an OTPK).  Star-Mesh handshake figures follow the $M_0$ definition of
Section~\ref{sec:x3dh}: $IK^{DSA}_{pk,A}$ (1{,}952~B) $+$ $IK^{DH}_{pk,A}$ (32~B) $+$
$EK_{pk,A}$ (32~B) $+$ $ct_1$ (1{,}088~B) $+$ optionally $ct_2$ (1{,}088~B) $+$
$\approx$32~B framing/identifiers.  The 1{,}952-byte $IK^{DSA}_{pk,A}$ is transmitted on
\emph{every} handshake as specified; implementations that cache a correspondent's identity
key after first contact can amortise this cost across a relationship rather than a session.}
\label{tab:overhead}
\begin{tabular}{@{}llll@{}}
\toprule
\textbf{Component}               & \textbf{Classical (Signal)} & \textbf{Star-Mesh} & \textbf{Overhead} \\
\midrule
Handshake, with OTPK (bytes)     & $\approx 150$ \cite{marlinspike_x3dh} & $\approx 4{,}224$ & $+4{,}074$ \\
Handshake, SPK-only fallback     & $\approx 118$ \cite{marlinspike_x3dh} & $\approx 3{,}136$ & $+3{,}018$ \\
Ratchet header (non-PQ step)     & 32 (X25519 pk)              & 32 (X25519 pk)     & $0$ \\
Ratchet header (PQ step)         & N/A                         & $+1{,}184$ (KEM pk) & $+1{,}184$ \\
Ratchet reply (PQ step)          & N/A                         & $+1{,}088$ (KEM ct) & $+1{,}088$ \\
\bottomrule
\end{tabular}
\end{table}

\paragraph{Amortised Cost.}
The PQ ratchet step adds $1{,}184 + 1{,}088 = 2{,}272$ bytes per round-trip.  Let $\Delta = 50$
be the ratchet interval and let $\delta_{\text{DHT}}$ denote the per-message DHT overhead
(key-routing records and mailbox polling headers).  The round-trip spans two independent
directions: the ML-KEM public key ($|PQ\_pk| = 1{,}184$~B) is carried in Alice's outgoing
messages amortised over $\Delta$ transmissions, and the corresponding ciphertext
($|PQ\_ct| = 1{,}088$~B) is returned in Bob's replies amortised over a further $\Delta$
receptions.  Computing the per-message overhead averaged across both directions gives a
divisor of $2\Delta$:\footnote{Equivalently, the per-message-in-one-direction cost is
$|PQ\_pk|/\Delta = 23.68$~B outgoing and $|PQ\_ct|/\Delta = 21.76$~B incoming; the
$2\Delta$ formulation is the symmetric average.}
\[
  \begin{aligned}
    \text{Per-message overhead}
      &= \frac{|PQ\_pk| + |PQ\_ct|}{2\Delta} + \delta_{\text{DHT}} \\
      &= \frac{1{,}184 + 1{,}088}{2 \times 50} + \delta_{\text{DHT}}
      = 22.72 + \delta_{\text{DHT}} \;\text{ bytes.}
  \end{aligned}
\]
In our Kademlia deployment, a mailbox routing record consists of a 20-byte node ID plus an
8-byte TTL field, and a polling header is 32 bytes.  As an illustrative engineering estimate
(the exact figure depends on routing topology, replication factor, churn model, and polling
cadence), amortised across 50 messages this contributes $\delta_{\text{DHT}} \approx 1$--$10$
bytes per message, giving a total of approximately $23.7$--$32.7$ bytes per message.  For a
standard 200-byte text message, this represents a modest overhead of approximately
$12\%$--$16\%$, becoming negligible for larger payloads or media transmissions.

\paragraph{Computational Cost.}
ML-KEM-768 encapsulation and decapsulation execute in well under 1\,ms on modern hardware
\cite{bos2018kyber}, triggered at most once every 50 messages.
To validate the computational practicality of our protocol, we micro-benchmarked our Rust
prototype over 1{,}000 iterations on an Apple M1 Air (8-core CPU, macOS 26.5.2), compiled in
Rust \texttt{release} mode (\texttt{opt-level = 3}, link-time optimisation enabled).  On this
platform, the complete hybrid PQ-X3DH handshake (comprising Alice's initiator computations,
Bob's responder computations, and the corresponding KEM encapsulation and decapsulation)
completes in an average of $17.93$~ms, and the post-quantum ratchet step (key generation,
encapsulation, decapsulation, and Root Key mixing) achieves a round-trip latency of $8.73$~ms.

These figures are higher than the sub-millisecond primitive benchmarks reported for Kyber in
\cite{bos2018kyber}, which measure optimised assembly implementations of the raw
encapsulation/decapsulation operations in isolation.  The gap is explained by the sequential
composition of multiple sub-operations in each handshake or ratchet step (two or three
ML-KEM-768 operations, two or four X25519 scalar multiplications, one HKDF-SHA3-256 derivation,
and the BLAKE3-KDF symmetric ratchet step), together with Rust runtime allocation, stack-frame
costs, and RustCrypto library overhead absent from the C/assembly benchmarks of \cite{bos2018kyber}.
A hardware-optimised, AVX2-enabled Rust build is expected to bring individual KEM operations
close to the cited sub-millisecond figures; the end-to-end handshake latency is ultimately
bounded by network round-trip time rather than CPU cost.
These execution times confirm that the computational overhead of hybrid post-quantum operations
is highly practical for standard consumer hardware.  The dominant computational overhead in
Star-Mesh is expected to be the mixnet onion routing (Sphinx packet construction \cite{sphinx}),
which requires multiple asymmetric operations per hop and is independent of the PQC primitives.

% ============================================================
\section{Assumptions and Limitations}
\label{sec:assumptionslimitations}
% ============================================================

We consolidate the protocol's assumptions and known limitations here for clarity.

\begin{enumerate}
  \item \textbf{Module-LWE hardness} (Assumption~\ref{asm:mlwe}): ML-KEM-768 is IND-CCA2
        secure under MLWE with $n = 256$, $k = 3$.
  \item \textbf{HKDF as PRF} (Assumption~\ref{asm:hkdf}): HKDF-SHA3-256 is modeled as a
        two-stage PRF following \cite{krawczyk2010hkdf}.
  \item \textbf{Deletion schedule enforced} (Assumption~\ref{asm:deletion}): An operational
        assumption; not proven by the protocol logic but required for FS and PCS.
  \item \textbf{RNG integrity after $t_{rec}$} (Assumption~\ref{asm:rng}): The RNG is clean
        from $t_{rec}$ onward.  A compromised RNG during PQ ratchet key generation makes PCS
        information-theoretically unachievable.
  \item \textbf{Mixnet honesty fraction} (Assumption~\ref{asm:mixnet}): At most $f < 1/2$ of
        mix nodes are adversarially controlled.
  \item \textbf{DHT Sybil resistance} (Assumption~\ref{asm:dht}): The DHT is operationally
        protected against Sybil and eclipse attacks; this is not proven by the protocol.
  \item \textbf{Out-of-order delivery and PCS} (Corollary~\ref{cor:pcs}): PCS restoration is
        conditional on the round-trip completing; delayed network delivery defers healing.
  \item \textbf{Skipped-key cache and PCS}: Keys cached in $\mathcal{K}_{skipped}$ before
        $t_{rec}$ remain compromised even after the PQ ratchet round-trip.
  \item \textbf{Implicit authentication only in $M_0$}: Star-Mesh provides only implicit
        authentication for $M_0$.  Explicit session-level mutual authentication (e.g., a
        signed transcript) is left for future work.
\end{enumerate}

% ============================================================
\section{Related Work}
\label{sec:related}
% ============================================================

\paragraph{Signal and PQXDH.}
Signal's PQXDH \cite{signal_pqxdh} is the current industry standard for post-quantum initial
key agreement.  It integrates Kyber into the X3DH handshake to provide PQ-FS against HNDL
attacks, similarly to our PQ-X3DH.  However, PQXDH (i)~relies on centralized servers for key
distribution and message routing, (ii)~does not integrate a post-quantum ratchet for ongoing
sessions, and (iii)~uses Ed25519 (classical) for identity signatures, leaving identity
authentication vulnerable to a CRQC.  Star-Mesh is designed to address all three gaps.
The fundamental technical obstacle PQXDH does not confront is serverless pre-key distribution:
moving post-quantum key bundles onto a decentralized DHT requires defending against eclipse
and key-substitution attacks on multi-kilobyte PQ pre-key material without any trusted
central arbiter, a problem that does not arise in Signal's centralized server model.

\paragraph{Session Protocol and Briar.}
Session \cite{session_protocol} routes messages through an onion network to achieve metadata
resistance and eliminates central servers.  Briar \cite{briar_design} routes through Tor.  Both
achieve excellent decentralization and metadata resistance but use exclusively classical X25519
for confidentiality, leaving them vulnerable to HNDL attacks.  Star-Mesh integrates hybrid PQC
into a similar decentralized architecture.
The technical obstacle these systems do not solve is the size incompatibility between
post-quantum primitives and onion-routing packet disciplines: Tor and Sphinx-based mixnets
enforce fixed-size cells or constant-size packets to prevent traffic analysis, but ML-DSA
signatures and ML-KEM ciphertexts exceed these bounds, requiring non-trivial fragmentation or
padding strategies that directly weaken the traffic-analysis resistance those networks are
designed to provide.

\paragraph{Double Ratchet Formal Analysis.}
Cohn-Gordon et al.\ \cite{cohngordon2016} provided the first formal security analysis of the
Signal Double Ratchet.  Cremers and Zhao\ \cite{dowling2020ratchet} extended this to a more
fine-grained model.  Our security model and argument structure are closely adapted from these
works, with aligned freshness/oracle semantics but protocol-specific reformulation for the
hybrid post-quantum setting with ML-KEM.
The technical obstacle these analyses leave open is the adaptation of ratchet freshness
semantics to KEM-based key exchange: classical DH-ratchet models assume commutativity and
ephemerality properties that do not hold for KEMs, where the encapsulator and decapsulator
roles are asymmetric and static public keys must be pre-distributed, forcing a new formal
treatment of RNG-compromise recovery and one-time pre-key exhaustion that is absent from
existing Double Ratchet proofs.

\paragraph{PQ-MLS and Group Messaging.}
RFC~9420 \cite{rfc9420} defines the base Messaging Layer Security (MLS) protocol for
centralized group messaging; post-quantum extensions to MLS are an active area of
standardization effort.  Star-Mesh targets asynchronous pairwise messaging over a DHT.
The two protocol classes are complementary and address different deployment scenarios.
The technical obstacle that prevents MLS-style designs from operating over a decentralized
DHT is their dependence on a central Delivery Service to impose a total ordering on group
operations and maintain tree consistency: without such a server, concurrent pre-key updates
produce conflicting tree states that cannot be resolved without a global coordination
mechanism, a problem that does not arise in Star-Mesh's pairwise ratchet setting.

\paragraph{Hybrid KEM Composition.}
Giacon et al.\ \cite{giacon2018kem} proved that the concatenation combiner for KEMs
achieves IND-CCA2 security if either component KEM is IND-CCA2 secure, provided the KEM ciphertexts are bound into the KDF input and the combiner
uses a suitable KDF modeled as a PRF (as established in \cite{giacon2018kem}).  Our
hybrid KEM construction satisfies this precondition by binding the ciphertexts $ct_1, ct_2$ and the ephemeral public key $EK_{pk,A}$ into the HKDF $info$ parameter, applying the composition results of \cite{giacon2018kem} with HKDF as
the PRF (Assumption~\ref{asm:hkdf}).
However, these static composition results do not extend to the stateful ratcheting setting:
they bound the advantage of a single IND-CCA2 distinguisher against a one-shot encapsulation,
but do not address the session-level forward secrecy and post-compromise recovery properties
that arise when KEM shared secrets are fed into a continuously updated ratchet state machine
where cached skip-keys and out-of-order delivery create additional exposure windows.

\paragraph{Metadata Resistance.}
Chaum \cite{chaum_mix} introduced mix networks; Piotrowska et al.\ \cite{loopix} developed the
Loopix architecture providing provable anonymity under formal assumptions.  The Sphinx packet
format \cite{sphinx} provides onion encryption with constant-size packets.  Star-Mesh inherits these results
directly for its routing layer.
The technical obstacle these systems leave unaddressed is the integration of large
post-quantum key material into a constant-packet-size discipline: Loopix's anonymity proof
assumes all packets are padded to a fixed size, but ML-DSA-87 signatures ($>4.5$~KB) and
ML-KEM-1024 ciphertexts ($1.5$~KB) far exceed standard Sphinx MTUs, and na\"ively introducing
these payloads forces fragmentation or re-padding that invalidates the constant-size assumption
on which the anonymity bound depends.  Star-Mesh resolves this by selecting Category~3
parameters (ML-KEM-768, ML-DSA-65) that keep the full authenticated pre-key bundle within
standard MTU bounds, as argued in Section~\ref{sec:paramselection}.

% ============================================================
\section{Conclusion}
\label{sec:conclusion}
% ============================================================

We presented Star-Mesh, a \emph{hybrid} post-quantum messaging protocol combining a
hybrid PQ-X3DH handshake (X25519 + ML-KEM-768) with a PQ Double Ratchet state machine,
Kademlia DHT key distribution, and mixnet-based metadata resistance.  We gave formal game-based
security definitions for PQ-FS and PCS in the quantum-active adversary model, with a unified
freshness definition that integrates the RevealRNG oracle.  We argued security via game-based
reductions showing that both properties follow from the IND-CCA2 security of ML-KEM-768 via
explicit named hybrid game sequences.  Authentication is similarly argued to follow from EUF-CMA
of ML-DSA-65.  The bundle signature covers all pre-key material (including OTPK and PQ-OTPK
batch hashes) to prevent substitution attacks.  The protocol's overhead is modest: amortised
per-message PQ overhead is approximately $22.7 + \delta_{\text{DHT}}$ bytes at $\Delta = 50$,
where $\delta_{\text{DHT}} \approx 1$--$10$ bytes is an illustrative engineering estimate
subject to routing topology and deployment configuration.  All assumptions are listed
explicitly in Section~\ref{sec:assumptions} and Section~\ref{sec:assumptionslimitations}.

Key assumptions underlying the security claims are: Module-LWE hardness
(Assumption~\ref{asm:mlwe}), the deletion schedule being enforced by the implementation
(Assumption~\ref{asm:deletion}), a clean RNG from $t_{rec}$ onward (Assumption~\ref{asm:rng}),
at most $f < 1/2$ honest mix nodes (Assumption~\ref{asm:mixnet}), and HKDF-SHA3-256 modeled
as a PRF (Assumption~\ref{asm:hkdf}).

\paragraph{Research Positioning.}
Star-Mesh should be viewed as a protocol-design and analysis contribution rather than a
production-ready messaging system.  The present manuscript focuses on architectural synthesis,
adversarial modeling, and hybrid post-quantum ratcheting semantics.  A production deployment
would additionally require extensive implementation hardening, empirical evaluation,
interoperability testing, formal verification, and operational abuse-resistance analysis.

\paragraph{Future Work.}
Several directions remain open.  First, a full composable security proof (e.g., in the UC
framework \cite{uc_framework}) for the complete protocol would be desirable.  Second, group
messaging extensions analogous to PQ-MLS \cite{rfc9420} should be explored for the DHT
setting.  Third, an implementation and empirical performance evaluation on resource-constrained
devices (IoT, mobile) would validate the practical viability of the design.  Fourth, an
explicit session-authentication mechanism for $M_0$ (to complement the current implicit
authentication) is a natural extension.

\section*{Acknowledgements}
The author thanks the anonymous reviewers for their time and feedback.  AI-assisted writing
tools were used solely for language refinement and typographical editing of this manuscript.
All protocol designs, security definitions, security arguments, and technical contributions are
entirely the work of the author.

% ============================================================
% Bibliography
% ============================================================
\begin{thebibliography}{99}

\bibitem{nist_fips203}
National Institute of Standards and Technology (NIST).
\emph{FIPS 203: Module-Lattice-Based Key-Encapsulation Mechanism Standard (ML-KEM)}.
Federal Information Processing Standard, August 2024.
\url{https://doi.org/10.6028/NIST.FIPS.203}

\bibitem{nist_fips204}
National Institute of Standards and Technology (NIST).
\emph{FIPS 204: Module-Lattice-Based Digital Signature Standard (ML-DSA)}.
Federal Information Processing Standard, August 2024.
\url{https://doi.org/10.6028/NIST.FIPS.204}

\bibitem{signal_pqxdh}
Signal Foundation.
\emph{The PQXDH Key Agreement Protocol}.
Technical Report, 2023.
\url{https://signal.org/docs/specifications/pqxdh/} (Accessed: July 2026).

\bibitem{marlinspike_x3dh}
M. Marlinspike and T. Perrin.
\emph{The X3DH Key Agreement Protocol}.
Signal Foundation, 2016.
\url{https://signal.org/docs/specifications/x3dh/} (Accessed: July 2026).

\bibitem{marlinspike_ratchet}
M. Marlinspike and T. Perrin.
\emph{The Double Ratchet Algorithm}.
Signal Foundation, 2016.
\url{https://signal.org/docs/specifications/doubleratchet/} (Accessed: July 2026).

\bibitem{cohngordon2016}
K. Cohn-Gordon, C. Cremers, B. Dowling, L. Garratt, and D. Stebila.
\emph{A Formal Security Analysis of the Signal Messaging Protocol}.
In \emph{Proceedings of the 2nd IEEE European Symposium on Security and Privacy (EuroS\&P 2017)},
pp.~451--466, IEEE, 2017.
\url{https://doi.org/10.1109/EuroSP.2017.27}

\bibitem{dowling2020ratchet}
C. Cremers and M. Zhao.
\emph{Secure Messaging with Strong Compromise Resilience, Temporal Privacy, and Immediate Decryption}.
Cryptology ePrint Archive, Report 2022/1481, 2022.
\url{https://eprint.iacr.org/2022/1481}

\bibitem{kademlia}
P. Maymounkov and D. Mazi\`eres.
\emph{Kademlia: A Peer-to-Peer Information System Based on the XOR Metric}.
In \emph{Proceedings of the 1st International Workshop on Peer-to-Peer Systems (IPTPS~2002)},
LNCS 2429, pp.~53--65, Springer, 2002.
\url{https://doi.org/10.1007/3-540-45748-8_5}

\bibitem{chaum_mix}
D. L. Chaum.
\emph{Untraceable Electronic Mail, Return Addresses, and Digital Pseudonyms}.
\emph{Communications of the ACM}, 24(2):84--90, 1981.
\url{https://doi.org/10.1145/358549.358563}

\bibitem{loopix}
A. M. Piotrowska, J. Hayes, T. Elahi, S. Meiser, and G. Danezis.
\emph{The Loopix Anonymity System}.
In \emph{Proceedings of the 26th USENIX Security Symposium}, pp.~1199--1216, USENIX, 2017.
\url{https://www.usenix.org/conference/usenixsecurity17/technical-sessions/presentation/piotrowska}

\bibitem{sphinx}
G. Danezis and I. Goldberg.
\emph{Sphinx: A Compact and Provably Secure Mix Format}.
In \emph{Proceedings of the 30th IEEE Symposium on Security and Privacy (S\&P~2009)},
pp.~269--282, IEEE, 2009.
\url{https://doi.org/10.1109/SP.2009.15}

\bibitem{lyubashevsky2010mlwe}
V. Lyubashevsky, C. Peikert, and O. Regev.
\emph{On Ideal Lattices and Learning with Errors over Rings}.
In \emph{Proceedings of EUROCRYPT 2010}, LNCS 6110, pp.~1--23, Springer, 2010.
\url{https://doi.org/10.1007/978-3-642-13190-5_1}

\bibitem{bos2018kyber}
J. Bos, L. Ducas, E. Kiltz, T. Lepoint, V. Lyubashevsky, J. Schanck, P. Schwabe, G. Seiler,
and D. Stehl\'e.
\emph{CRYSTALS\,--\,Kyber: A CCA-Secure Module-Lattice-Based KEM}.
In \emph{Proceedings of the 3rd IEEE European Symposium on Security and Privacy (EuroS\&P~2018)},
pp.~353--367, IEEE, 2018.
\url{https://doi.org/10.1109/EuroSP.2018.00032}

\bibitem{giacon2018kem}
F. Giacon, F. Heuer, and B. Poettering.
\emph{KEM Combiners}.
In \emph{Proceedings of PKC 2018}, LNCS 10769, pp.~190--218, Springer, 2018.
\url{https://doi.org/10.1007/978-3-319-76578-5_7}

\bibitem{bindel2018hybrid}
N. Bindel, J. Brendel, M. Fischlin, B. Goncalves, and D. Stebila.
\emph{Hybrid Key Encapsulation Mechanisms and Authenticated Key Exchange}.
In \emph{Proceedings of PQCrypto 2019}, LNCS 11505, pp.~206--226, Springer, 2019.
\url{https://doi.org/10.1007/978-3-030-25510-7_12}

\bibitem{blake1997unknown}
S. Blake-Wilson and A. Menezes.
\emph{Unknown Key-Share Attacks on the Station-to-Station (STS) Protocol}.
In \emph{Proceedings of PKC~1999}, LNCS 1560, pp.~154--170, Springer, 1999.
\url{https://doi.org/10.1007/3-540-49162-7_12}

\bibitem{mosca2018cybersecurity}
M. Mosca.
\emph{Cybersecurity in an Era with Quantum Computers: Will We Be Ready?}
\emph{IEEE Security \& Privacy}, 16(5):38--41, 2018.
\url{https://doi.org/10.1109/MSP.2018.3761723}

\bibitem{session_protocol}
Session Foundation.
\emph{Session: Private Messenger Technical Overview}.
\url{https://docs.getsession.org/}, 2023. (Accessed: July 2026).

\bibitem{briar_design}
The Briar Project.
\emph{Briar Design Documentation}.
\url{https://briarproject.org/}, 2023. (Accessed: July 2026).

\bibitem{rfc9420}
R. Barnes, B. Beurdouche, R. Robert, J. Millican, E. Omara, and K. Cohn-Gordon.
\emph{The Messaging Layer Security (MLS) Protocol}.
IETF RFC 9420, July 2023.
\url{https://www.rfc-editor.org/rfc/rfc9420}

\bibitem{rfc7748}
A. Langley, M. Hamburg, and S. Turner.
\emph{Elliptic Curves for Security}.
IETF RFC 7748, January 2016.
\url{https://www.rfc-editor.org/rfc/rfc7748}

\bibitem{rfc5869}
H. Krawczyk and P. Eronen.
\emph{HMAC-based Extract-and-Expand Key Derivation Function (HKDF)}.
IETF RFC 5869, May 2010.
\url{https://www.rfc-editor.org/rfc/rfc5869}

\bibitem{krawczyk2010hkdf}
H. Krawczyk.
\emph{Cryptographic Extraction and Key Derivation: The HKDF Scheme}.
In \emph{Proceedings of CRYPTO 2010}, LNCS 6223, pp.~631--648, Springer, 2010.
\url{https://doi.org/10.1007/978-3-642-14623-7_34}

\bibitem{blake3}
J. O'Connor, J.-P. Aumasson, S. Neves, and Z. Wilcox-O'Hearn.
\emph{BLAKE3: One Function, Fast Everywhere}.
Technical Report, 2020.
\url{https://github.com/BLAKE3-team/BLAKE3-specs} (Accessed: July 2026).

\bibitem{nist_aesgcm}
National Institute of Standards and Technology (NIST).
\emph{Recommendation for Block Cipher Modes of Operation: Galois/Counter Mode (GCM) and GMAC}.
NIST Special Publication 800-38D, November 2007.
\url{https://doi.org/10.6028/NIST.SP.800-38D}

\bibitem{uc_framework}
R. Canetti.
\emph{Universally Composable Security: A New Paradigm for Cryptographic Protocols}.
In \emph{Proceedings of the 42nd IEEE Symposium on Foundations of Computer Science (FOCS~2001)},
pp.~136--145, IEEE, 2001.
\url{https://doi.org/10.1109/SFCS.2001.959888}

\bibitem{dodis2013entropic}
Y. Dodis, T. Ristenpart, and S. Vadhan.
\emph{Randomness Condensers for Efficiently Samplable, Seed-Dependent Sources}.
In \emph{Proceedings of TCC 2012}, LNCS 7194, pp.~618--635, Springer, 2012.
\url{https://doi.org/10.1007/978-3-642-28914-9_35}

\end{thebibliography}

\end{document}
